\SourceLeaf{321}
\[
q=-\cos\chi\mathord{.}x-\tfrac12\operatorname{tang}Q\mathord{.}\sin\chi^2xx\quad .\ .\ .
\]
vornimmt (ihre Herleitung ist einfach und kann hier entfallen) und das Ergebnis mit der Reihe
\[
m=1+\mu x^3+\mu'x^4\quad .\ .\ .
\]
vergleicht. Für unseren gegenwärtigen Zweck genügt es jedoch zu zeigen, dass die gesuchte Länge des geodätischen Bogens von \(Ah\) um höchstens eine Größe fünfter Ordnung abweicht. Da \(h^5+10h^3\delta\delta+5h\delta^4\) ersichtlich eine solche Größe ist, braucht der entwickelte Wert von \(\mu'\) hier nicht angegeben zu werden. Für \(\mu\) ergibt sich dagegen der Wert
\[
\mu=\frac{2ee\cos P\mathord{.}\sin P}{3\cos\varphi\cos\theta}\mathord{.}\cos\chi^3
\]
und da \(\delta\cos\chi\) nach Artikel 15 eine Größe zweiter Ordnung ist, ist offenbar auch \((h^3\delta+h\delta^3)\mu\) eine Größe fünfter Ordnung.

Da \(h\) dasselbe bedeutet, was jetzt mit \(s\) bezeichnet wird, gilt demnach bis zur fünften Ordnung ausschließlich genau
\[
s=\frac rA\quad .\ .\ .\ .\ .\ .\ .\ .\ .\ .\ .\ .\ .\ .\ .\quad(8)
\]

\indent Damit schließlich alles für die weitere Entwicklung Nötige an einer Stelle steht, gebe ich hier noch die folgenden exakten Gleichungen an, die schon in der ersten Abhandlung (Artikel 4, 6 und 3) verwendet wurden,
\[
\begin{aligned}
A&=\frac{a\cos\varphi}{\cos\theta^2}\quad .\ .\ .\ .\ .\ .\ .\ .\ .\ .\ .\ .\ .\ .\quad(9)\\
\cos Q&=\frac{\cos\theta\cos P}{\alpha\cos\varphi}\quad .\ .\ .\ .\ .\ .\ .\ .\ .\ .\ .\ .\ .\quad(10)\\
\sin Q&=\frac{\sin P}{\alpha}
\end{aligned}
\]
sowie die Gleichung, die sich aus der Verbindung der beiden letzten ergibt
\[
\operatorname{tang}Q=\frac{\cos\varphi\operatorname{tang}P}{\cos\theta}\quad .\ .\ .\ .\ .\ .\ .\ .\ .\ .\ .\ .\quad(11)
\]

\Article{23}
Für unseren Zweck brauchen diese Gleichungen jetzt nur noch passend kombiniert zu werden.
